\chapter{基于改进YOLO11-Seg的立体仓库堆叠货物盘点方法}{Method for Stacked Cargo Inventory Based on Improved YOLO11-Seg}

本章围绕立体仓库堆叠货物盘点任务中的检测、分割与计数问题展开研究。首先分析基线模型 YOLO11s-Seg 的结构特点及其在复杂堆叠场景中的局限；其次在此基础上提出改进模型 YOLO11-StackLite，并对其关键模块进行设计说明；最后结合实例分割结果与空间几何关系，构建后处理数量计算方法，为后续实验验证与模型压缩研究奠定基础。

\section{堆叠货物检测与分割}{Stacked Cargo Detection and Segmentation}

\subsection{YOLO11s-Seg基础架构}{}
为实现立体仓库中堆叠货物的高效自动盘点，视觉模型不仅需要准确区分堆叠结构中的顶层箱体与非顶层箱体，还需要在边缘部署条件下具备较低的计算开销与较快的推理速度。因此，本文选取兼顾检测效率与实例分割能力的单阶段模型作为研究基础。

基于上述需求，本文选择 Ultralytics 发布的 YOLO11s-Seg 作为基线模型。该模型在单阶段框架下集成了目标检测与实例分割能力，在检测精度与推理速度之间取得了较好平衡，具备较强的边缘部署适配性。YOLO11s-Seg 的整体结构主要由主干网络（Backbone）、特征融合网络（Neck）和检测头（Head）构成。

在 YOLO11s-Seg 的 Backbone 与 Neck 中，引入了改进的 C3k2 Fast 模块（CSP Bottleneck with Two Convolution Layers），用于增强浅层特征提取与融合能力。该模块基于跨阶段部分结构思想，能够在保持较低计算复杂度的同时提升特征表达能力。在 Backbone 末端，集成轻量级注意力模块 C2PSA，用于增强模型对关键区域的感知能力，其结构如\autoref{fig:fig2}所示。

\begin{figure}[H]
    \centering
    \includegraphics[width=0.8\textwidth]{fig2}
    \bicaption{YOLO11s-Seg主干与注意力结构示意}{Backbone and attention structure of YOLO11s-Seg}
    \label{fig:fig2}
\end{figure}

C2PSA 采用 Softmax Attention，其计算公式如下：
\begin{equation}\label{eq:attn}
  \mathrm{Attention}(Q,K,V)=\mathrm{softmax}\left(\frac{QK^{\mathsf T}}{\sqrt{d_k}}\right)V
\end{equation}
其中，$Q$、$K$、$V$ 分别表示查询、键和值。该注意力机制通过计算不同位置之间的相关性实现全局信息建模，其时间复杂度为 $O(N^2)$。尽管该结构能够增强模型对重要区域的关注能力，但在高分辨率特征图上，其计算代价仍然较高。此外，原始注意力机制在复杂堆叠场景中对低相关区域的抑制能力有限，容易引入冗余信息，从而影响对关键结构特征的有效表征。

因此，虽然 YOLO11s-Seg 已具备较好的基础性能，但在复杂堆叠场景下仍存在注意力计算冗余、下采样效率不足以及空间结构建模能力有限等问题，有必要在其基础上进行针对性的结构改进。

\subsection{YOLO11-StackLite架构}{}
鉴于仓储盘点系统对识别精度、推理效率与部署资源均具有较高要求，尽管 YOLO11s-Seg 在常规场景下已表现出较好的综合性能，但在复杂堆叠图像中仍可能出现漏检、误检等问题，且其参数规模与计算复杂度在边缘部署场景下仍存在进一步优化空间。为此，本文在 YOLO11s-Seg 的基础上提出改进模型 YOLO11-StackLite，从注意力建模、下采样效率和空间结构表达三个方面对网络进行优化，以提升复杂堆叠场景下的检测与分割性能，并进一步降低模型部署开销。

改进后的实例分割网络模型称为 YOLO11-StackLite，其整体结构如\autoref{fig:fig3}所示。
\begin{figure}[H]
    \centering
    \includegraphics[width=0.8\textwidth]{fig3}
    \bicaption{YOLO11-StackLite整体结构}{Overall structure of YOLO11-StackLite}
    \label{fig:fig3}
\end{figure}

围绕上述目标，本文设计了 3 个关键改进模块，即通道拆分注意力模块 C2SASA、重排分组卷积模块 SSConv，以及具有四向卷积结构的 C3k2\_QDConv 模块。三者分别面向关键区域建模、轻量化下采样和方向性空间特征提取，从不同层面增强模型对堆叠货物场景的适应能力。

\par\noindent 1）通道拆分注意力模块 C2SASA：该模块在继承原始注意力结构优势的基础上，通过通道拆分机制与选择性激活抑制策略，降低注意力计算冗余并增强模型对关键区域的表征能力。

\par\noindent 2）重排分组卷积模块 SSConv：该模块面向 Neck 部分的轻量化下采样需求，通过融合像素重排、分支卷积与组卷积等设计，在降低计算复杂度的同时提升多尺度特征提取能力。

\par\noindent 3）具有四向卷积的 C3k2\_QDConv 模块：该模块通过增强对方向性几何结构的感知能力，改善传统 Bottleneck 在复杂结构建模中的不足，从而提高模型在堆叠场景下的检测与分割性能。

\subsubsection{通道拆分注意力模块 C2SASA}{}
为同时提升模型对显著区域的建模能力并降低计算开销，本文提出通道拆分式选择性激活抑制注意力模块（C2SASA），用于替换 Backbone 中的 C2PSA 模块。该模块首先通过 $1\times1$ 卷积对输入特征进行通道压缩，并将其划分为两条分支：其中一条分支保留残差信息以维持原始特征表达，另一条分支输入改进注意力块进行特征增强，最后通过通道拼接与卷积映射恢复原始维度，实现两类信息的有效融合。

在改进注意力块中，本文将原始 Attention 替换为 SASA 机制，以实现对特征的选择性激活与抑制。SASA 模块在保持全局建模能力的同时，引入 L2 归一化与可学习温度参数，从而增强注意力计算过程中的数值稳定性和特征区分能力，其结构如\autoref{fig:fig4}所示。

\begin{figure}[H]
    \centering
    \includegraphics[width=0.75\textwidth]{fig4}
    \bicaption{选择性激活抑制注意力机制 SASA}{Selective activation suppression attention}
    \label{fig:fig4}
\end{figure}

设输入张量为 $x\in R^{B\times C\times H\times W}$，其中 $B$ 为批大小，$H$、$W$ 分别表示输入特征图的高度和宽度，$C$ 为特征维度。首先将空间维度展平并转置为序列形式，其中 $N=H\cdot W$。随后采用共享线性投影层一次性生成 $q$、$k$、$v$：
\begin{equation}\label{eq:qkv}
[q,k,v] = {\rm{Linear}}_{qkv}(X)
\end{equation}
为实现多头注意力，将通道维度 $C$ 重构为 $h\times d$，其中 $h$ 为注意力头数，$d=C/h$ 为每个注意力头的嵌入维度。由此得到：
\[
q,k,v \to R^{B \times h \times N \times d}
\]
对 $q$ 与 $k$ 进行 L2 归一化：
\[
q' = \frac{q}{{\left\| q \right\|_2}},\quad
k' = \frac{k}{{\left\| k \right\|_2}}
\]
进一步引入可学习温度参数 $\mathrm{Temp}\in\mathbb{R}^{h\times1\times1}$，得到相似度矩阵：
\begin{equation}\label{eq:sim}
\mathrm{sim} = \mathrm{Temp} \cdot q'k'^{\mathsf T}
\end{equation}
其中，$\mathsf T$ 表示转置操作。随后通过 Softmax 与 Dropout 获得注意力权重：
\[
\alpha = {\rm{Dropout}}\left( {{\rm{softmax}}\left( {{\rm{sim}}} \right)} \right)
\]
最后对 $v$ 进行加权求和并投影：
\[
y = {\rm{Proj}}\left( {\alpha  \cdot v} \right)
\]
最终输出通过 Transpose 与 Reshape 操作恢复到原始空间维度。综上，C2SASA 在保留全局依赖建模能力的基础上，降低了注意力计算的冗余性，并增强了对关键区域的特征区分能力，为复杂堆叠场景下的实例检测与分割提供了更有效的特征表示。

\subsubsection{重排分组卷积模块 SSConv}{}
在 YOLO11 模型中，Neck 部分通常采用串联的标准卷积进行下采样。虽然该结构实现简单，但存在参数量较大、计算开销较高以及通道信息利用不足等问题，不利于模型在边缘设备上的高效运行。与此同时，标准卷积在下采样过程中对局部细节信息和多尺度语义信息的保留能力有限，容易影响复杂堆叠场景中的特征表达效果。

针对上述不足，本文提出重排分组卷积模块 SSConv，以实现下采样阶段的轻量化设计与特征增强。该模块融合了像素重排、分支卷积、组卷积与逐点卷积等结构思想，通过通道重组和双分支处理策略，在降低模型计算成本的同时增强特征提取能力，其结构如\autoref{fig:fig5}所示。

\begin{figure}[H]
    \centering
    \includegraphics[width=0.75\textwidth]{fig5}
    \bicaption{重排分组卷积模块 SSConv}{Shuffle-split convolution}
    \label{fig:fig5}
\end{figure}

该模块用于替换 YOLO11-Seg 中 Neck 部分的标准下采样卷积结构，以提高下采样阶段的效率和特征利用能力。设输入特征图为 $X$，首先通过 PixelShuffle 操作增强通道间信息交互，再沿通道维度将特征划分为两个子分支 $X_1$ 和 $X_2$：
\[
{X_1},{X_2} = {\rm{Split}}\left( {{\rm{PixelShuffle}}\left( X \right),\dim  = 1} \right)
\]
第一条分支采用组卷积与逐点卷积进行轻量化特征提取与跨通道融合：
\[
{Y_1} = {\rm{GroupConv}}\left( {{{\rm{X}}_1}} \right)+{\rm{PointwiseConv}}\left( {{{\rm{X}}_1}} \right)
\]
第二条分支采用步幅为 2 的 $3\times3$ 卷积，以保持较大的感受野并增强对整体结构的感知能力：
\[
{Y_2} = {\rm{Conv2}}{{\rm{d}}_{3 \times 3}}\left( {{{\rm{X}}_2}} \right)
\]
最终将两路特征进行拼接并经过归一化与激活：
\[
Y = {\rm{SiLU}}\left( {{\rm{BN}}\left( {{\rm{Concat}}\left( {{Y_1}{\rm{,}}{Y_2}} \right)} \right)} \right)
\]
因此，SSConv 在降低下采样阶段计算负担的同时，兼顾了局部细节保持与跨通道信息融合能力，有助于提升模型在复杂堆叠场景中的多尺度特征表达效果。

\subsubsection{具有四向卷积的 C3k2\_QDConv 模块}{}
在堆叠货物场景中，箱体边缘、遮挡关系以及方向性几何结构对目标识别具有重要影响。然而，传统 Bottleneck 模块主要依赖标准卷积完成局部特征提取，其对方向性结构和复杂空间关系的建模能力相对有限。尤其在库存边界模糊、局部重叠遮挡等情况下，传统卷积往往难以充分刻画长条形边缘和多方向变化特征，容易造成细节信息损失。

为增强模型对堆叠场景几何特征的感知能力，本文提出 Bottleneck\_QDConv，并构建对应的 C3k2\_QDConv 模块。该模块以四向卷积为核心，通过多个方向上的非对称卷积与显式空间位移操作，加强对不同方向边缘和空间结构变化的建模能力，其结构如\autoref{fig:fig6}所示。

\begin{figure}[H]
    \centering
    \includegraphics[width=0.8\textwidth]{fig6}
    \bicaption{四向卷积模块 QDConv}{Quad-direction convolution module}
    \label{fig:fig6}
\end{figure}

设输入特征为 $X$，输出为 $Y$，卷积核尺寸为 $k$，位移步长为 $s=1$。首先对输入特征进行零填充，得到四个方向的位移分支 $X_p^1$、$X_p^2$、$X_p^3$ 和 $X_p^4$。随后分别对其施加 $1\times k$ 与 $k\times 1$ 的卷积，并通过显式位移操作进行方向建模，最终表达式为：
\[
Y = \mathrm{Conv}_{2\times 2}\left(
\mathrm{Concat}\left[
\begin{array}{l}
\mathrm{Roll}_{+1}^{H}\bigl(\mathrm{Conv}_{1\times k}(X_p^1)\bigr),\\
\mathrm{Roll}_{-1}^{H}\bigl(\mathrm{Conv}_{1\times k}(X_p^2)\bigr),\\
\mathrm{Roll}_{+1}^{W}\bigl(\mathrm{Conv}_{k\times 1}(X_p^3)\bigr),\\
\mathrm{Roll}_{-1}^{W}\bigl(\mathrm{Conv}_{k\times 1}(X_p^4)\bigr)
\end{array}
\right]
\right)
\]
其中，$\mathrm{Roll}_{+1}^{H}$、$\mathrm{Roll}_{-1}^{H}$ 表示沿高度方向进行正向和反向位移操作；$\mathrm{Roll}_{+1}^{W}$、$\mathrm{Roll}_{-1}^{W}$ 表示沿宽度方向进行正向和反向位移操作；$\mathrm{Conv}_{1\times k}$、$\mathrm{Conv}_{k\times 1}$ 分别表示核尺寸为 $1\times k$ 和 $k\times 1$ 的卷积操作，用于水平方向和垂直方向的特征提取。

与传统标准卷积相比，QDConv 不仅能够在多个方向上捕获目标的几何变化特征，还能够通过位移操作强化空间关系表达，从而使模型在复杂堆叠结构下获得更丰富的方向性特征表示。由此，C3k2\_QDConv 能够提升模型对复杂空间结构和目标轮廓信息的刻画能力，为后续检测与分割提供更具判别性的特征支撑。

\section{后处理数量计算方法}{Post-processing Counting Method}
在立体仓库堆叠货物盘点任务中，仅依赖目标检测结果难以准确完成数量估计，其原因在于盘点结果不仅取决于目标是否被识别，还与上下层结构关系及堆叠规则密切相关。由于实例分割结果能够提供更精细的轮廓信息和空间形状信息，本文进一步结合实例分割掩膜设计后处理数量计算方法，以实现对复杂堆叠场景下货物数量的准确推断。

\subsection{计数主体流程}{}
本文采用“up\_box”和“down\_box”两个类别对每幅图像中的目标进行标注，如\autoref{fig:fig7}所示。
\begin{figure}
  \subfloat[原图]{\label{subfig:fig7a}
    \includegraphics[width=0.4\textwidth]{fig7-1.png}
  }
  \qquad
  \subfloat[标注后]{\label{subfig:fig7b}
    \includegraphics[width=0.4\textwidth]{fig7-2.png}
  }
  \bicaption{图像标注示例}{Example of image annotation}
  \label{fig:fig7}
\end{figure}
其中，类别标签 1 和 0 分别对应上层货物与下层货物。模型在推理阶段输出每个目标实例的边界框、类别标签及其对应的掩膜信息。为了提升后续计数环节的准确性与鲁棒性，本文根据不同类别的特性设置差异化的置信度阈值，仅保留满足阈值要求的预测结果参与后续处理。

在获得满足阈值条件的实例分割结果后，本文设计了一种基于实例分割结果的启发式后处理计数方法，其总体流程如\autoref{fig:fig8}所示。根据仓库货物“下一层箱体仅在当前层堆满后方可叠加”的堆叠规则，本文将总货物数量估算划分为下层箱体数量估算与顶层箱体数量推断两个阶段。
\begin{figure}[H]
    \centering
    \includegraphics[width=0.8\textwidth]{fig8}
    \bicaption{后计数算法流程图}{Flow chart of post counting algorithm}
    \label{fig:fig8}
\end{figure}

根据预测结果中上下层目标的存在状态，本文将堆叠货物盘点情形划分为以下三类：

a）检测结果为空。若预测结果中未检测到任何目标，即类别数组为空，则可直接判定该图像中无货物，数量记为 0。

b）仅检测到下层货物。若预测结果中仅包含类别 0，而不包含类别 1，则说明图像中缺乏顶层结构信息，无法直接判断完整堆叠关系，如\autoref{subfig:fig9a}所示。此时输出 \texttt{num=-1}，并将该类图像交由仓储管理系统进一步处理。

c）同时检测到上层与下层目标。该情形为最常见的完整堆叠场景，需要进一步区分顶层满载和顶层未满载两种情况。若顶层判定为满载，则总货物数量可依据上下层目标数量关系直接估算：
\begin{equation}\label{eq:total-full}
{N_{total}} = {N_{up\_box}} \times ({N_{down\_box}} + 1)
\end{equation}
式中，$N_{total}$ 为货物总数量，$N_{up\_box}$ 与 $N_{down\_box}$ 分别为上层和下层目标数量。若顶层未满载，则需进一步结合掩膜关键点与几何关系推理完整堆叠结构，并估算每层下层箱体数量 $N_{layer}$，最终总货物数量表示为：
\begin{equation}\label{eq:total-unfull}
{N_{total}} = {N_{up\_box}} + ({N_{layer}} \times {N_{down\_box}})
\end{equation}

\begin{figure}
  \subfloat[缺乏顶层信息]{\label{subfig:fig9a}
    \includegraphics[width=0.25\textwidth]{9-1.png}
  }
  \qquad
  \subfloat[顶层满载]{\label{subfig:fig9b}
    \includegraphics[width=0.23\textwidth]{9-2.png}
  }
  \qquad
  \subfloat[顶层未满载]{\label{subfig:fig9c}
    \includegraphics[width=0.24\textwidth]{9-3.png}
  }
  \bicaption{立体仓库堆叠货物盘点的三种情况}{Three situations for inventory of stacked goods in a three-dimensional warehouse}
  \label{fig:fig9}
\end{figure}

下面将分别对顶层满载检测算法以及基于掩膜关键点的层间面积比估计算法进行介绍，以说明本文后处理数量计算方法的具体实现过程。

\subsubsection{基于几何关系的顶层满载检测算法}{}
当场景中同时存在上下层目标时，首先需要判断顶层是否处于满载状态。为此，本文提出一种基于掩模几何关系的顶层满载检测算法，其流程见表~\ref{tab:alg_full_load}。该算法首先对实例掩模进行轮廓提取，并利用凸包运算获取边界角点，以近似描述货物的平面外形与空间占位关系。随后，根据模型预测标签将实例划分为上层货物集合与下层货物集合，并依据图像坐标系下的几何位置关系比较二者在垂直方向上的相对位置，从而判断上层目标是否覆盖下层顶部区域。

这种基于几何中心与边界点的比较方式，能够在无深度信息的条件下利用二维平面几何关系实现有效的层间遮挡判断。因此，该算法能够作为后续数量估计的前置判别步骤，为不同盘点情形下的数量推断提供依据。

\begin{table}[H]
  \bicaption{基于几何关系的顶层满载检测算法}{Top-layer full-load detection algorithm based on geometric relationships}
  \label{tab:alg_full_load}
  \centering
  \begin{tabular}{cl}
    \toprule
    步骤 & 具体内容 \\
    \midrule
    输入 & 预测标签列表 \texttt{pred\_labels}，掩模列表 \texttt{masks} \\
    输出 & 是否满载（布尔值） \\
    1 & 初始化角点集合 \texttt{pts\_list} \\
    2 & 遍历所有实例掩模并提取轮廓 \\
    3 & 对有效轮廓计算凸包角点 \texttt{pts} \\
    4 & 根据预测标签划分上层与下层目标集合 \\
    5 & 计算下层目标最高点的纵坐标 \texttt{dn\_top\_y} \\
    6 & 计算上层目标凸包中心的纵坐标 \texttt{up\_center\_y} \\
    7 & 若 \texttt{up\_center\_y > dn\_top\_y}，则判定为满载 \\
    8 & 否则判定为未满载 \\
    \bottomrule
  \end{tabular}
\end{table}

\subsubsection{掩膜关键点驱动的层间面积比估计算法}{}
当顶层未满载时，无法直接依据满层假设完成总数量估计，因此需要进一步估算其下方次顶层箱体数量。针对这一问题，本文提出一种基于掩膜关键点驱动的层间面积比估计方法，其流程见表~\ref{tab:alg_area_ratio}，关键点示意见\autoref{fig:fig10}。

在计算过程中，需要求解不同候选边之间的交点，以确定上下层结构之间的空间重叠关系。设任意两条线段分别为 $P_1P_2$ 与 $Q_1Q_2$，其中 $P_1(x_A,y_A)$、$P_2(x_B,y_B)$、$Q_1(x_C,y_C)$、$Q_2(x_D,y_D)$，其斜率可表示为：
\begin{equation}\label{eq:line-k}
{k_1} = \frac{{{y_B} - {y_A}}}{{{x_B} - {x_A}}}, \quad
{k_2} = \frac{{{y_D} - {y_C}}}{{{x_D} - {x_C}}}
\end{equation}
根据两条直线的一般式方程，可进一步写为矩阵形式：
\begin{equation}\label{eq:line-intersection}
\left( \begin{array}{l}
{a_1}{\rm{  }}{b_1}\\
{a_2}{\rm{  }}{b_2}
\end{array} \right)\left( \begin{array}{l}
x\\
y
\end{array} \right) = \left( \begin{array}{l}
 - {c_1}\\
 - {c_2}
\end{array} \right)
\end{equation}
其中，$(x,y)$ 为两条线段的交点坐标。公式\autoref{eq:line-k}和\autoref{eq:line-intersection}用于求解候选边对的交点集合，并为后续面积比估计提供几何基础。

\begin{table}[H]
  \bicaption{掩膜关键点驱动的层间面积比算法}{Interlayer area-ratio algorithm driven by mask keypoints}
  \label{tab:alg_area_ratio}
  \centering
  \begin{tabular}{cl}
    \toprule
    步骤 & 具体内容 \\
    \midrule
    输入 & 上层掩膜集合 \texttt{up\_box\_masks}，下层掩膜集合 \texttt{down\_box\_masks} \\
    输出 & 估算的次顶层箱体数量 \texttt{Nlayer} \\
    1 & 提取次顶层关键点 $A$、$B$、$C$、$D$、$G$ \\
    2 & 根据 $AB$ 与 $CD$ 的几何关系求解交点 $O_1$ \\
    3 & 提取顶层关键点 $\_A$、$\_B$、$\_C$、$\_D$、$\_E$、$\_F$ \\
    4 & 对候选边对求交，得到顶层交点集合，并取平均位置作为 $O_2$ \\
    5 & 计算顶层单箱体区域面积 $S_1$ \\
    6 & 计算次顶层整体区域面积 $S_2$ \\
    7 & 由面积比得到 $N_{layer}=S_2/S_1$ \\
    8 & 返回估计结果 \texttt{Nlayer} \\
    \bottomrule
  \end{tabular}
\end{table}

上述方法充分利用了实例分割掩膜所提供的轮廓与几何信息，能够在顶层未满载的情况下提升数量估计的可解释性和准确性，从而增强盘点算法对复杂堆叠场景的适应能力。

\begin{figure}[H]
    \centering
    \includegraphics[width=0.8\textwidth]{fig10}
    \bicaption{目标掩膜凸包点和关键点示意图}{Schematic diagram of convex hull points and key points on the target mask}
    \label{fig:fig10}
\end{figure}

本章首先分析了基线模型 YOLO11s-Seg 的结构特点及其在复杂堆叠场景中的局限性，在此基础上提出改进模型 YOLO11-StackLite，并从注意力建模、下采样结构和空间方向感知三个方面对关键模块进行了设计。随后，结合实例分割结果与空间几何关系，构建了面向堆叠货物场景的后处理数量计算方法。上述研究为提升复杂场景下的检测、分割与盘点性能提供了方法基础。然而，为满足边缘部署需求，仍需进一步降低模型的参数量与计算复杂度，因此下一章将围绕模型压缩方法展开研究。
